\section{Research gaps}
\label{related-work:research-gaps}
Three gaps in the literature surveyed above motivate the research questions of this thesis.

First, network transfer accounting is absent from almost all spatial benchmarks, despite being a primary driver of cost and latency in cloud delivery. \textcite{tang2020_locationspark_distributed_spatial_query} treats network communication as a system-design concern rather than a reported experimental metric, and \textcite{folkerts2012_benchmarking_in_the_cloud} argues that cloud benchmarks must report end-to-end performance jointly with pricing, of which data transfer is a major component. The link from format internals to client-observed cost has therefore not been measured directly in the spatial benchmarking literature.

Second, no peer-reviewed benchmark compares vector cloud-native formats (GeoParquet, PMTiles) against traditional formats and engines (Shapefile, PostGIS, WMS-based vector tiles) on the same workload. The empirical comparisons that exist appear in community and vendor reports \parencite{holmes2023_geoparquet_performance, flatgeobuf2024_flatgeobuf_benchmarks}, are restricted to local-disk timing and file size, and do not place cloud-native and traditional formats side by side on cloud storage.

Third, no published benchmark compares single-node engines (PostGIS, DuckDB) against distributed engines (Apache Sedona on Databricks) on the same vector workload. \textcite{pandey2018_how_good_modern_spatial_analytics} compares only Spark-based systems, while Jackpine \parencite{ray2011_jackpine_spatial_database_benchmark} and similar single-node benchmarks predate the Spark-based spatial systems that emerged in the mid-2010s.

The first two gaps motivate \acrshort{rq}1, which compares cloud-native formats and engines against traditional ones in performance and operational cost across analytical and visualization spatial workloads. The third gap motivates \acrshort{rq}2, which extends the comparison to horizontal scaling of distributed cloud-native engines against single-node engines on the same workloads. The methodological gap discussed in Section~\ref{ch:methodology}, the absence of effect-size reporting and uncertainty quantification in geospatial benchmarks, is addressed as a methodological contribution rather than a separate research question.